\honest{Just Lose Hope Already}{Eliezer Yudkowsky}{2007}

Casey Serin, a 24-year-old web programmer with no experience in real estate, owes banks \$2.2 million. He lied on mortgage applications to buy eight houses in different states at once, borrowed more than the houses cost, and spent the difference on living expenses and real-estate seminars. He was expecting the market to go up, it seems. \nb{Contemporary reports, among them ABC's \textsc{Nightline} in April 2007, confirm his age, his job, the eight houses and the \$2.2 million of debt.}

That is not the sad part. The sad part is that he has not given up. He will not declare bankruptcy or get a job, he keeps buying seminars, and he tried to take out a mortgage on a ninth house. He has not failed, you see; he has had a learning experience. ``That's what happens when you refuse to lose hope.''

This puts me in mind of two Nobel-prize-winning economists, Merton and Scholes of Long-Term Capital Management. \nb{Merton and Scholes were partners in the fund, which was founded and run by John Meriwether. A Federal Reserve history reports that Scholes was the partner most uneasy about the fund's move into new markets in 1997.} After three years of huge profits, the tricks the fund exploited stopped working, because others had learned them. The fund refused to lose hope. Addicted to 40\% annual returns, it borrowed more and more to exploit smaller and smaller margins. When things started to go wrong, it had \$4.72 billion of equity, \$124.5 billion of borrowing and \$1.25 trillion in derivatives. \nb{In the crisis of 1998, the same history reports, Meriwether ``still believed that their trades would pay off if they could simply weather the storm.'' Wikipedia reports that the positions, taken over by a group of banks, were eventually sold at a small profit to them. The post does not describe the collapse.}

Every profession has its own skills, so you might think a general study of rationality would add little. But it seems to me that not being stupid, not turning little mistakes into big ones, is much the same art ``whether in hedge funds or romance.'' One of its keys: ``Be ready to admit you lost.'' I give no example from romance. \nb{The post gives no way to tell a loss from a setback. If the account of the rescue is right, LTCM's trades were not simply lost, and what sank the fund was, on this reading, the borrowing, which left it unable to wait.}
